\setcounter{section}{3} % Continuing from the previous segment

% ==========================================
% SECTION 4: COMPUTING THE UNIT DISK
% ==========================================
\section{Computing the Area of a Unit Disk}

\begin{spoken_clean}[00:15:00 - 00:18:37]
Now that we have established the rigorous analytical definition of $\pi$, we will prove—perhaps for the first time in your mathematical lives—that the area of the unit disk is exactly $\pi$. 

In our rigorous language, we are asking for the 2-dimensional Jordan measure of the set $B_1$, which is the set of all points $(x_1, x_2)$ such that $x_1^2 + x_2^2 \le 1$. To compute this measure, we will use polar coordinates. 

Let the first parameter $y_1$ represent the radius, and the second parameter $y_2$ represent the angle. We define our coordinate mapping $\Psi$ such that $x_1 = y_1 \cos(y_2)$ and $x_2 = y_1 \sin(y_2)$. Our parameter domain $D$ will be the open rectangle where the radius $y_1$ goes from $0$ to $1$, and the angle $y_2$ goes from $0$ to $2\pi$.
\end{spoken_clean}

\begin{nice-box}[Setting up the Transformation]
The professor writes down the set definition for the disk and formally defines the polar mapping $\Psi$ alongside its parameter domain $D$.
\end{nice-box}

\begin{math_stroke}[Formalizing the Polar Transformation]
The Unit Disk in the original coordinates:
\[
B_1 = \{ (x_1, x_2) \in \mathbb{R}^2 \mid x_1^2 + x_2^2 \le 1 \}
\]
The Polar Coordinate Mapping $\Psi$:
\begin{align*}
x_1 &= y_1 \cos(y_2) \\
x_2 &= y_1 \sin(y_2)
\end{align*}
The Parameter Domain $D$:
\[
D = \underbrace{(0,1)}_{\text{Radius } y_1} \times \underbrace{(0, 2\pi)}_{\text{Angle } y_2} \implies \Psi: D \to B_1
\]

\textbf{Explanation of Steps:}
We are transitioning from Cartesian coordinates $(x_1, x_2)$ to polar coordinates $(y_1, y_2)$. The domain $D$ is strictly an \textit{open} rectangle. This strict open boundary guarantees that $\Psi$ remains a proper diffeomorphism (bijective and smooth), as including $y_1 = 0$ (the origin) or $y_2 = 0, 2\pi$ (the overlapping seam) would break the bijection.
\end{math_stroke}

\begin{spoken_clean}[00:18:37 - 00:20:17]
Our first crucial observation is that the image of our parameter domain, $\Psi(D)$, is not entirely equal to the closed disk $B_1$. Because we are using the open interval $(0,1)$ for the radius and $(0, 2\pi)$ for the angle, our mapping completely misses the origin and the line segment lying on the positive x-axis. 

However, we can easily prove that this missing 1-dimensional line segment has a 2-dimensional Jordan measure of zero. If you cover it with small squares of side length $\epsilon$, the total area goes to zero. Therefore, even though the sets are not identical, their measures are perfectly equal: $\mu_2(B_1) = \mu_2(\Psi(D))$.
\end{spoken_clean}

\begin{didactic_insight}[The "Measure Zero" Safety Net]
This is a critical Real Analysis concept. The fact that a mapping isn't perfectly surjective over the boundaries introduces the "cheat code" of integration: integrals do not "see" sets of measure zero. Boundaries, seams, and single points can be freely deleted without changing the overall volume.
\end{didactic_insight}

\begin{nice-box}[TikZ: The Measure Zero Boundary Issue]
The professor draws the mapping, actively highlighting the "missing" seam on the right side of the circle.
\end{nice-box}

\begin{math_stroke}[The Measure Zero Boundary Issue]
\begin{center}
\begin{tikzpicture}[scale=1.5]
    % Domain D
    \draw[thick, fill=profblue!10] (0,0) rectangle (1, 2);
    \node at (0.5, 1) {$D$};
    \node[below] at (0,0) {$0$};
    \node[below] at (1,0) {$1$};
    \node[left] at (0,2) {$2\pi$};
    \node[below] at (0.5, -0.5) {Parameter Space $(y_1, y_2)$};

    % Mapping Arrow
    \draw[->, very thick] (1.5, 1) -- (2.5, 1) node[midway, above] {$\Psi$};

    % Disk B1
    \draw[thick, fill=profblue!10] (4,1) circle (1cm);
    \draw[dashed, profred, very thick] (4,1) -- (5,1); % The missing seam
    
    % Annotations
    \node[profred, right] at (5,1) {\textbf{Missing Seam}};
    \node[profred, right] at (5,0.7) {$\mu_2 = 0$};
    \node at (4, 1.5) {$B_1$};
    \node[below] at (4, -0.5) {Target Space $(x_1, x_2)$};
\end{tikzpicture}
\end{center}
\end{math_stroke}

\begin{spoken_clean}[00:20:17 - 00:22:55]
Because the measures are equal, we can apply our Change of Variables substitution rule. The area of the disk is the integral of the constant function $1$ over our domain $D$, multiplied by the determinant of the Jacobian matrix of $\Psi$. 

Let us compute this Jacobian matrix. Taking the partial derivatives with respect to $y_1$ and $y_2$, we construct our $2 \times 2$ matrix. When we compute the determinant, we get $y_1 \cos^2(y_2) + y_1 \sin^2(y_2)$. Factoring out the $y_1$, we use our analytical axiom that $\cos^2 + \sin^2 = 1$ (i.e., $y_1(\cos^2(y_2) + \sin^2(y_2)) = y_1$). The determinant is simply $y_1$. Because $y_1$ is strictly positive on our domain, its absolute value is just $y_1$ (i.e., $|\det J\Psi| = y_1$).
\end{spoken_clean}


\begin{math_stroke}[Calculating the Jacobian Determinant]
\textbf{Applying the Substitution Rule:}
\[
\mu_2(B_1) = \int_{\Psi(D)} 1 \, dx = \int_D 1 \cdot \underbrace{|\det J\Psi(y_1, y_2)|}_{\text{Distortion Factor}} \, dy
\]
\textbf{Computing the Jacobian Matrix:}
\[
J\Psi(y_1, y_2) = \begin{pmatrix} 
\frac{\partial x_1}{\partial y_1} & \frac{\partial x_1}{\partial y_2} \\[6pt]
\frac{\partial x_2}{\partial y_1} & \frac{\partial x_2}{\partial y_2}
\end{pmatrix} = \begin{pmatrix} 
\cos(y_2) & -y_1 \sin(y_2) \\ 
\sin(y_2) & y_1 \cos(y_2) 
\end{pmatrix}
\]
\textbf{Evaluating the Determinant:}
\begin{align*}
\det J\Psi &= \big(\cos(y_2) \cdot y_1 \cos(y_2)\big) - \big(-y_1 \sin(y_2) \cdot \sin(y_2)\big) \\
&= y_1 \cos^2(y_2) + y_1 \sin^2(y_2) \\
&= y_1 \underbrace{(\cos^2(y_2) + \sin^2(y_2))}_{=1 \text{ (Pythagorean identity)}} \\
&= y_1
\end{align*}

\textbf{Explanation of Steps:}
The Jacobian determinant tells us exactly how much a tiny square of parameter space $dy_1 dy_2$ is stretched when it is mapped into the disk. Because the determinant is $y_1$, it means the "stretching" scales linearly with the radius. Grids near the outer edge of the circle are stretched much larger than grids near the origin.
\end{math_stroke}

\begin{spoken_clean}[00:22:55 - 00:26:33]
Substituting this back in, we must compute the double integral of $y_1$ over our rectangular domain $D$ (i.e., $\int_D y_1 \, dy$). But we have a problem: we haven't learned how to systematically compute multivariable integrals yet! 

We can solve this using a geometric trick. The integral of a positive function over a 2D domain represents the 3-dimensional volume of its \textit{hypograph}—the physical region trapped underneath the function's surface. 

Let's add a vertical coordinate, $x_3$. We want the volume where $x_3 \le y_1$, over our base rectangle $(0,1) \times (0, 2\pi)$. 
Consider a full 3D bounding box with dimensions $1 \times 2\pi \times 1$. The total volume is $2\pi$. Our function, $x_3 = y_1$, forms a slanted plane that cuts this box perfectly in half along its diagonal. Therefore, the volume of our hypograph is exactly half the volume of the box, which gives us $\pi$ (i.e., $\frac{1}{2} (1 \cdot 2\pi \cdot 1) = \pi$).
\end{spoken_clean}

\begin{nice-box}[The Hypograph Wedge]
The professor draws a 3D bounding box and shows how the function $f(y_1, y_2) = y_1$ cuts it precisely in half.
\end{nice-box}

\begin{math_stroke}[Evaluating the Volume via Geometry]
\begin{equation} \label{eq:disk_integral}
\mu_2(B_1) = \int_{(0,1) \times (0, 2\pi)} \underbrace{y_1}_{\text{Height } x_3} \, \underbrace{dy_1 \, dy_2}_{\text{Base Area}}
\end{equation}

\begin{center}
\begin{tikzpicture}[scale=1.5]
    % Bounding Box Coordinates (1 x 2pi x 1)
    \coordinate (O) at (0,0,0);
    \coordinate (X) at (4,0,0); % Length 2pi along y2
    \coordinate (Y) at (0,2,0); % Height 1 along x3
    \coordinate (Z) at (0,0,2); % Depth 1 along y1
    \coordinate (XY) at (4,2,0);
    \coordinate (XZ) at (4,0,2);
    \coordinate (YZ) at (0,2,2);
    \coordinate (XYZ) at (4,2,2);

    % Back dashed lines
    \draw[dashed, gray] (O) -- (X);
    \draw[dashed, gray] (O) -- (Y);
    \draw[dashed, gray] (O) -- (Z);

    % The Hypograph Wedge (x3 <= y1)
    % The plane is x3 = y1. At y1=0 (front), height is 0. At y1=1 (back), height is 1.
    \filldraw[fill=profred!30, draw=profred, thick, opacity=0.8] (X) -- (XZ) -- (XYZ) -- cycle; % Right side triangle
    \filldraw[fill=profred!20, draw=profred, thick, opacity=0.8] (O) -- (X) -- (XYZ) -- (YZ) -- cycle; % Slanted roof x3=y1
    \filldraw[fill=profred!40, draw=profred, thick, opacity=0.8] (O) -- (Z) -- (YZ) -- cycle; % Left side triangle
    
    % Remaining box outline
    \draw[thick, gray] (Z) -- (XZ) -- (XYZ) -- (YZ) -- cycle; 
    \draw[thick, gray] (X) -- (XY) -- (XYZ);
    \draw[thick, gray] (Y) -- (XY);
    \draw[thick, gray] (Y) -- (YZ);

    % Axis Labels
    \node[below] at (2,0,0) {Longitude $y_2 \in [0, 2\pi]$};
    \node[left] at (0,0,1) {Radius $y_1 \in [0, 1]$};
    \node[left] at (0,1,0) {Height $x_3$};
    \node[profred, right] at (4.2, 1, 1) {Hypograph Volume ($x_3 \le y_1$)};
\end{tikzpicture}
\end{center}

\textbf{Explanation of Steps:}
By interpreting the integral physically, we see the function $y_1$ creates a wedge. Because the wedge is a perfect diagonal bisection of a rectangular prism, we can compute its volume using simple geometry instead of advanced integration rules.
\[
\mu_2(B_1) = \frac{1}{2} \Vol(\text{Bounding Box}) = \frac{1}{2} \big(1 \cdot 2\pi \cdot 1\big) = \pi
\]
\end{math_stroke}

% ==========================================
% SECTION 5: VOLUME OF THE 3D UNIT BALL
% ==========================================
\section{Example: Volume of the Unit Ball in $\mathbb{R}^3$}

\begin{spoken_clean}[00:26:33 - 00:32:00]
That trick is wonderful, but let us move to a harder question. Let's compute the volume of the unit ball in $\mathbb{R}^3$. We know from geometry that the answer should be $\frac{4}{3}\pi$, but we have not proved it analytically yet.

We will try the exact same approach using 3D polar coordinates—spherical coordinates. We like spherical coordinates because they flawlessly transform rectangular boxes into spheres. Let our mapping be $\Psi(y_1, y_2, y_3)$. 

The first coordinate, $y_1$, is always the radius. The second coordinate, $y_2$, is the longitude—the deviation from a fixed meridian. The third coordinate, $y_3$, represents the elevation. Let me write the formulas down...
\end{spoken_clean}

\begin{nice-box}[The Coordinate Convention Trap]
The professor writes the spherical coordinates. He intends for $y_3$ to be the "elevation" (latitude, measured from the equator). However, he writes the standard physics convention where the angle is measured from the North Pole (colatitude). This is mathematically incorrect for his stated domain.
\end{nice-box}

\begin{math_stroke}[The Initial Incorrect Formulation]
\begin{align} 
x_1 &= y_1 \cos(y_2) \sin(y_3) \nonumber \\
x_2 &= y_1 \sin(y_2) \sin(y_3) \label{eq:wrong_spherical} \\
x_3 &= y_1 \cos(y_3) \nonumber
\end{align}
\end{math_stroke}

\begin{didactic_insight}[The Coordinate Trap]
There are two competing standards for spherical coordinates. Mathematics often uses latitude (angle from the equator, meaning $z = r\sin\phi$). Physics usually uses colatitude (angle from the North Pole, meaning $z = r\cos\phi$). Mixing the words of one convention with the math of the other inadvertently creates a phenomenal active-learning moment for the room.
\end{didactic_insight}

\begin{spoken_clean}[00:32:00 - 00:32:53]
Wait, looking at this... did I make a mistake?
\end{spoken_clean}

\begin{meta_note}[Student Spotting the Error]
\textit{"If we would like $y_3$ to represent the elevation from the equator, shouldn't we switch the sine and cosine in the formula for the polar coordinates?"}
\end{meta_note}

\begin{spoken_clean}[continued]
Ah! The $y_3$ is the elevation, right? Yes, you are completely right. If we use the cosine for the $x_3$ coordinate, we are measuring the deviation from the North Pole. If we want the elevation from the equator, the vertical $x_3$ coordinate must be governed by the sine function. Let me fix this on the board. The $x_3$ gets the sine, and the others get the cosine. Excellent catch.
\end{spoken_clean}

\begin{nice-box}[Correcting the Spherical Coordinates]
The professor erases the sines and cosines and formally swaps them.
\end{nice-box}

\begin{math_stroke}[The Corrected Spherical Coordinates]
The mapping $\Psi: D \to \mathbb{R}^3$ for elevation-based spherical coordinates:
\begin{align} 
x_1 &= y_1 \cos(y_2) \underbrace{\cos(y_3)}_{\text{Corrected}} \nonumber \\
x_2 &= y_1 \sin(y_2) \underbrace{\cos(y_3)}_{\text{Corrected}} \label{eq:correct_spherical} \\
x_3 &= y_1 \underbrace{\sin(y_3)}_{\text{Corrected}} \nonumber
\end{align}

\textbf{Explanation of Steps:}
By setting $x_3 = y_1 \sin(y_3)$, we ensure that when the elevation angle $y_3 = 0$, our height $x_3$ is $0$ (putting us on the equator). When $y_3 = \pi/2$, $\sin(\pi/2) = 1$, putting us exactly at the North Pole $x_3 = y_1$.

The parameter domain $D$ is formally defined as:
\[
D = \underbrace{(0,1)}_{y_1 (\text{Radius})} \times \underbrace{(0, 2\pi)}_{y_2 (\text{Longitude})} \times \underbrace{\left(-\frac{\pi}{2}, \frac{\pi}{2}\right)}_{y_3 (\text{Elevation})}
\]
\end{math_stroke}

\begin{nice-box}[Spherical Coordinates Visualized]
Visual representation of the corrected coordinate system, explicitly showing $y_3$ sweeping up from the equatorial plane.
\end{nice-box}

\begin{math_stroke}[Spherical Coordinates Visualized]
\begin{center}
\begin{tikzpicture}[scale=2]
    % 3D Axes
    \draw[->, thick] (0,0) -- (-1.2,-1.2) node[below left] {$x_1$};
    \draw[->, thick] (0,0) -- (2.5,0) node[right] {$x_2$};
    \draw[->, thick] (0,0) -- (0,2.5) node[above] {$x_3$};

    % Sphere outline & Equator
    \draw[thick, gray!70] (0,0) circle (2);
    \draw[dashed, gray!70] (0,0) ellipse (2 and 0.6);
    \draw[gray!70] (-2,0) arc (180:360:2 and 0.6);

    % Point P and Projection
    \coordinate (O) at (0,0);
    \coordinate (P) at (1.2, 1.3);
    \coordinate (Pproj) at (1.2, -0.25); % Projection on equator

    % Construction Lines
    \draw[dashed, thick] (O) -- (Pproj);
    \draw[dashed, thick] (P) -- (Pproj);
    
    % Radius Vector
    \draw[thick, profblue, ->] (O) -- (P) node[midway, above left] {$y_1$ (Radius)};

    % Angles
    \draw[->, profred, thick] (-0.5,-0.5) to[bend right=20] node[below right, text=black] {$y_2$ (Longitude)} (0.5, -0.1);
    \draw[->, profgreen, thick] (0.6, -0.12) to[bend left=20] node[right, text=black] {$y_3$ (Elevation)} (0.6, 0.65);

    % Point Dot
    \fill[black] (P) circle (1.5pt) node[above right] {$(x_1, x_2, x_3)$};
\end{tikzpicture}
\end{center}
\end{math_stroke}
